% Materiał z AF 25; fragment do rozdziału o funkcjach testowych.
\subsection{Porównanie topologii przestrzeni funkcji testowych}
\label{sec:synteza25}

\begin{definicja}[Notacja porównawcza przestrzeni funkcji]
\label{def:trzy-przestrzenie25}
Zachowujemy przestrzenie $\mathcal E(\Omega)=C^\infty(\Omega)$,
$\mathcal D_K(\Omega)$, $\mathcal D(\Omega)$ i $\mathcal S(\R^d)$
z definicji \ref{def:DK18}, \ref{def:D-LF18}
oraz \ref{def:schwartz17}. Na $\mathcal E$ i $\mathcal D_K$
używamy seminorm
$p_{K,m}(f)=\max_{|\alpha|\le m}\sup_{x\in K}|\partial^\alpha f(x)|$.
Na $\mathcal S$ zapis $p_{N,m}(f)$ oznacza normę
$Q_{N,m}(f)$ z definicji \ref{def:schwartz17}.
Przestrzenie sprzężone topologicznie oznaczamy przez
$\mathcal E'$, $\mathcal D'$ i $\mathcal S'$; indeksy
$\sigma(X',X)$ i $\beta(X',X)$ oznaczają odpowiednio
topologię słabą-* (zbieżność punktowa na $X$)
i silną (zbieżność jednostajna na zbiorach ograniczonych).
Wszystkie przestrzenie rozpatrujemy nad wspólnym ciałem $\K$.
\end{definicja}

\begin{twierdzenie}[Ciągłe i gęste zanurzenia]
\label{thm:zanurzenia25}
Niech $d\ge1$. Odwzorowania identycznościowe
\[
 \mathcal D(\R^d)\longrightarrow\mathcal S(\R^d)
 \longrightarrow\mathcal E(\R^d)
\]
są ciągłymi, różnowartościowymi odwzorowaniami liniowymi o gęstych
obrazach. Dla dowolnego niepustego zbioru otwartego $\Omega\subseteq\R^d$
takie samo stwierdzenie zachodzi dla zanurzenia
$\mathcal D(\Omega)\to\mathcal E(\Omega)$.
Żadne z trzech wskazanych odwzorowań nie jest
izomorfizmem liniowo-topologicznym na swój obraz z topologią
podprzestrzeni; ostatnie stwierdzenie zachodzi dla każdego
niepustego zbioru otwartego $\Omega$.
\end{twierdzenie}
\begin{proof}
Dla ustalonego zwartego $K\Subset\R^d$ oraz $N,m\in\N_0$
mamy $p_{N,m}(f)\le (1+\sup_{x\in K}|x|)^N p_{K,m}(f)$
dla $f\in\mathcal D_K(\R^d)$. Stąd ciągłość
$\mathcal D_K\to\mathcal S$; własność uniwersalna granicy
induktywnej daje ciągłość $\mathcal D\to\mathcal S$.
Ciągłość $\mathcal S\to\mathcal E$ wynika z oszacowania
$p_{K,m}(f)\le p_{0,m}(f)$. Analogicznie sprawdzamy
ciągłość $\mathcal D(\Omega)\to\mathcal E(\Omega)$.

Aby wykazać gęstość, wybierzmy $\chi\in\mathcal D(\R^d)$
równe $1$ na kuli jednostkowej i zdefiniujmy
$\chi_R(x)=\chi(x/R)$ dla $R\ge1$. Dla $f\in\mathcal S(\R^d)$ i $|\alpha|\le m$
wzór Leibniza daje
\[
 \partial^\alpha((1-\chi_R)f)
 =\sum_{\gamma\le\alpha}\binom{\alpha}{\gamma}
 \partial^\gamma(1-\chi_R)\partial^{\alpha-\gamma}f.
\]
Wszystkie czynniki $\partial^\gamma(1-\chi_R)$
znikają dla $|x|<R$. Dla $\gamma\ne0$ są ograniczone
przez $R^{-|\gamma|}\sup|\partial^\gamma\chi|$,
a dla $\gamma=0$ przez $1+\sup|\chi|$.
Ponieważ $R\ge1$ i $(1+|x|)^{-1}\le(1+R)^{-1}$
na pozostałej części przestrzeni, otrzymujemy
\[
 p_{N,m}((1-\chi_R)f)
 \le \frac{C_{\chi,m}}{1+R}p_{N+1,m}(f).
\]
Zatem $\chi_Rf\to f$ w $\mathcal S$ i $\chi_Rf\in\mathcal D$.
Dla $f\in\mathcal E(\R^d)$ mamy $\chi_Rf=f$ w otoczeniu
każdego ustalonego zbioru zwartego przy dostatecznie dużym $R$,
a więc $\chi_Rf\to f$ w $\mathcal E$. W ogólnej
$\Omega$ używamy funkcji odcięcia $\chi_j\in\mathcal D(\Omega)$,
równej $1$ w otoczeniu $K_j$, gdzie $(K_j)$ jest
wyczerpaniem $\Omega$ zbiorami zwartymi. Każdy zwarty $L\Subset\Omega$ mieści się w pewnym $K_J$,
więc dla $j\ge J$ wszystkie pochodne różnicy znikają na $L$.
Stąd $\chi_j f\to f$ w $\mathcal E(\Omega)$.

Na $\R^d$ niech $\psi\in\mathcal D(\R^d)$ będzie niezerowa
i $f_n(x)=\psi(x-ne_1)$. Wówczas $f_n\to0$ w
$\mathcal E(\R^d)$, bo nośniki ostatecznie omijają każdy
ustalony zbiór zwarty, ale nie w $\mathcal S(\R^d)$,
gdyż $p_{0,0}(f_n)=p_{0,0}(\psi)>0$.
Z kolei $g_n(x)=e^{-n^2}\psi(x-ne_1)$ zbiega do zera w
$\mathcal S(\R^d)$: dla $\supp\psi\subseteq\overline B(0,A)$
mamy
\[
 p_{N,m}(g_n)\le e^{-n^2}(1+n+A)^N p_{0,m}(\psi)
 \longrightarrow0.
\]
Nie zbiega do zera w $\mathcal D(\R^d)$, gdyż nośniki
nie mieszczą się w jednym zbiorze zwartym
(Twierdzenie~\ref{thm:D-sequences18}).

Dla dowolnego niepustego otwartego $\Omega\subseteq\R^d$
wybieramy wyczerpanie $(K_j)$ jak w definicji.
Zbiór $\Omega\setminus K_j$ jest niepusty i otwarty,
bo żaden niepusty otwarty podzbiór $\R^d$ nie jest zwarty.
Wybieramy niezerowe $h_j\in\mathcal D(\Omega)$ takie,
że $\supp h_j\subseteq\Omega\setminus K_j$.
Wówczas $h_j\to0$ w $\mathcal E(\Omega)$, gdyż dla
każdego zwartego $L\Subset\Omega$ nośniki ostatecznie
omijają jego otoczenie (wybieramy $L\subseteq\operatorname{int}K_J$).
Gdyby $h_j\to0$ w $\mathcal D(\Omega)$, wszystkie nośniki
byłyby zawarte w pewnym zwartym $L\Subset\Omega$,
a więc w pewnym $K_J$. Dla $j\ge J$ oznaczałoby to
$h_j=0$, sprzecznie z wyborem. Topologia $\mathcal D(\Omega)$
jest zatem ściśle silniejsza od odziedziczonej z $\mathcal E(\Omega)$.
\end{proof}
